\section*{Градуировка и рёберный коцикл}

Цепь \(A_3\) имеет единственную \(\mathbb Z_2\)-градуировку с нечётными
рёбрами с точностью до общего знака: \(\Gamma=\operatorname{diag}(1,-1,1)\).
Но \(A_3\) является деревом, поэтому \(H^1(A_3;U(1))=0\): все фазы одной
общей связности снимаются вершинной калибровкой и дают нулевую CP.

Добавление хорды \(0\!\leftrightarrow\!2\) создаёт единственный поток
\(\Phi\). Для разбиения ``хорда против двухрёберной цепи'' точный CP-нечётный
инвариант пропорционален \(\sin\Phi\). Однако треугольник не двудолен, поэтому
все три ребра не могут одновременно быть нечётной частью одного
градуированного оператора Дирака. Минимальный положительный plaquette-action
также выбирает \(\Phi=0\), то есть сохраняет CP.